\section{Limitations}
\label{sec:limits}

Our results should be read as evidence for context curation as a design
principle, not as a universal dominance claim for any single policy. The
number of retained runs per cell is small, so we avoid significance claims
on individual cells and report run means, standard deviations, and
best-by-test operating points instead. The file-level strategy also shows
benchmark- and proposer-dependent trade-offs: it often produces lower-cost
operating points, but it can sacrifice passrate when the uncurated file pool
already contains strong signal or when guidance metadata is added on top of a
large exposed history.

The ablations are necessarily narrower than the main evaluation. The
Random@$3$ and Recency@$3$ iteration-selection ablations are anchored on the
LoCoMo / codex54 setting, while the force-budget analyses focus on retained
claudekimi runs. Likewise, the read-target mechanism analysis is most detailed
for claudekimi because its tool traces expose per-file read targets; codex54
logs bucket more activity under a single shell interface, limiting comparable
file-level attribution. Finally, \name is evaluated on three agent-harness
families, but all experiments use a fixed outer-loop length and a fixed set of
proposer backbones. Larger run budgets, additional proposers, and broader
task families may shift the cost--quality frontier and should be tested before
using the observed trade-offs as deployment defaults.
